\chapter{Extended Black-Box Complexity Theory}
\label{app:bbc-extended}

\section{Las Vegas vs.\ Monte Carlo Black-Box Complexity}
\label{app:lv-mc}

In the standard (unrestricted, unbiased) black-box models, the
\textit{Las Vegas} complexity — the expected number of queries until
an optimal solution is first found — and the \textit{Monte Carlo}
complexity — the minimum budget $T$ such that the optimum is found
with probability at least $1 - p$ for a given $p > 0$ — are
essentially equivalent up to constant factors via Markov's inequality:
any Las Vegas algorithm with expected runtime $T$ finds the optimum
within $T/p$ queries with probability at least $1 - p$.

In elitist models, however, this equivalence breaks down severely.
\citet{doerr2015onemax} show that the $(1+1)$ elitist Las Vegas and
Monte Carlo complexities of \textsc{OneMax} can be exponentially far
apart. The reason is structural: an elitist algorithm cannot perform
restarts, because a fresh random starting point will almost surely
have lower fitness than the current incumbent and will therefore be
rejected. Without restarts, a run that fails to find the optimum is
trapped, and there is no mechanism to recover. Consequently, failure
probabilities do not decay geometrically with additional budget in the
same way as for non-elitist algorithms.

Formally, the $p$-\textit{Monte Carlo black-box complexity} of a
problem class $\mathcal{F}$ with respect to an algorithm class
$\mathcal{A}$ is
\[
  \min_{A \in \mathcal{A}}\;
  \min\bigl\{T \in \mathbb{N} :
    \forall f \in \mathcal{F},\;
  \mathbb{P}(\tau_{f^*} \leq T) \geq 1 - p\bigr\}.
\]
The Las Vegas complexity is recovered in the limit by Markov's
inequality when the two notions coincide, but for the $(1+1)$ elitist
Las Vegas complexity of \textsc{OneMax}, \citeauthor{doerr2015onemax}
establish only that it is $O(n \log n)$ (matching the unary unbiased
lower bound) while leaving the exact constant open; the Monte Carlo
complexity is $\Theta(n)$, a separation of a full $\log n$
factor~\cite{doerr2015onemax}.

For the algorithms studied in this project, this distinction does not
affect the experimental results directly: all runtime and regret
analyses use expected quantities over independent runs, which
corresponds to the Las Vegas framework. The Monte Carlo / Las Vegas
distinction becomes relevant if one wishes to derive high-probability
tail bounds on cumulative regret, a direction noted as future work in
Section~\ref{subsec:open-questions}.

\section{Yao's Principle in Restricted Models}
\label{app:yao-restricted}

Yao's Minimax principle in \ref{subsubsec:yaomini} applies directly only to
algorithm classes whose convex hull contains the randomized algorithm
of interest. For memory-restricted and elitist models this containment
fails because some randomized algorithms in these classes are not
convex combinations of deterministic ones. The canonical example is
\acrshort{rls} itself: every deterministic $(1+1)$ memory-restricted
ranking-based algorithm that ever rejects a search point will enter an
infinite loop on \textsc{OneMax} with positive probability when the
input is chosen uniformly at random, giving infinite expected runtime,
whereas \acrshort{rls} optimizes \textsc{OneMax} in $\Theta(n \log n)$
expected queries~\cite{doerr2015onemax}.

\citet{doerr2015onemax} resolve this by applying Yao's principle to
the strictly larger class $\mathcal{A}'$ of all comparison-based
$(1+1)$ algorithms with access to their complete search history.
Every memory-restricted ranking-based algorithm can be simulated in
$\mathcal{A}'$, so a lower bound on the $\mathcal{A}'$-complexity
implies the same lower bound for the original class. The information-
theoretic argument then proceeds as in the unrestricted case: after
$i - 1$ queries the oracle has revealed at most one bit of information
per query, so at most $2^{i-1}$ candidate optima have been eliminated,
giving
\[
  \mathbb{P}(\text{optimum found by query } i) \leq 2^{i - n},
\]
and summing yields $\mathbb{E}[\tau_{f^*}] \geq n - 1$.
